% Terjemahan Bahasa Indonesia dari chapter2.tex baris 519--720.
% Sumber: Methods of Algebra, Volume 2, commit
% 9a5803ff77dd3257484cb177f851a73770a59dd3 (CC BY 4.0).
% stable-unit-id: o014.aljabr2.chapter2.glimpse-of-lattice-theory

% segment-id: o014.li2.chapter2.glimpse-of-lattice-theory.h001
\section{Sekilas tentang Teori Kisi}\label{sec:lattices}
\hypertarget{o014.aljabr2.chapter2.glimpse-of-lattice-theory}{ }

% segment-id: o014.li2.chapter2.glimpse-of-lattice-theory.p001
Bagian ini memberikan pengantar singkat tentang suatu kelas struktur terurut
parsial yang disebut kisi. Pembahasan ini akan membantu menjernihkan struktur
kategori abelian.

% segment-id: o014.li2.chapter2.glimpse-of-lattice-theory.p002
Sepanjang bagian ini, kita akan meninjau berbagai himpunan terurut parsial
$(P, \leq)$; kategori yang bersesuaian dengannya dilambangkan dengan
$\mathcal{P}$.

% segment-id: o014.li2.chapter2.glimpse-of-lattice-theory.p003
Untuk setiap himpunan bagian $S$ dari suatu himpunan terurut parsial
$(P, \leq)$, terdapat konsep batas atas, batas bawah, supremum $\sup S$, dan
infimum $\inf S$. Himpunan bagian $S \subset P$ mempunyai supremum (atau
infimum) jika dan hanya jika koproduk (atau produk) semua objek dalam $S$
terdapat di kategori $\mathcal{P}$; dalam hal ini supremum (atau infimum)
tersebut adalah koproduk (atau produk) itu.

% segment-id: o014.li2.chapter2.glimpse-of-lattice-theory.d001
\begin{definisi}\label{def:lattice-order}
	\index{kisi (lattice)}
	\index{kisi!terbatas (bounded)}
	\index{interval (interval)}
	\index[sym1]{veewedge@$\vee$, $\wedge$}
	Jika setiap dua unsur $a,b$ dari himpunan terurut parsial $(P, \leq)$
	mempunyai supremum $a \vee b$ dan infimum $a \wedge b$, maka $(P, \leq)$
	disebut \emph{kisi}. Jika, sebagai tambahan, $(P, \leq)$ sendiri mempunyai
	batas atas dan batas bawah, keduanya tunggal dan masing-masing dilambangkan
	dengan $1$ dan $0$; dalam hal ini $P$ disebut \emph{kisi terbatas}.

	Jika $a,b$ merupakan unsur himpunan terurut parsial $(P, \leq)$ dan
	$a \leq b$, maka $[a, b] := \left\{x \in P: a \leq x \leq b \right\}$ juga
	merupakan himpunan terurut parsial terhadap $\leq$ dan disebut
	\emph{interval} yang ditentukan oleh $a,b$.
\end{definisi}

% segment-id: o014.li2.chapter2.glimpse-of-lattice-theory.p004
Jika $P$ merupakan kisi, setiap interval dalam $P$ tertutup terhadap $\vee$
dan $\wedge$, serta merupakan kisi terbatas. Jika $P$ merupakan kisi terbatas,
maka $P=[0,1]$, dan untuk setiap $x \in P$ berlaku
$x \wedge 0=0$, $x \vee 0=x=x \wedge 1$, serta $x \vee 1=1$.

% segment-id: o014.li2.chapter2.glimpse-of-lattice-theory.p005
Dari sudut pandang kategori, $a \wedge b$, $a \vee b$, $1$, dan $0$ masing-masing
bersesuaian dengan produk dua objek, koproduk dua objek, objek terminal (yakni
produk kosong), dan objek awal (yakni koproduk kosong) di dalam kategori;
silakan pembaca memeriksanya.

% segment-id: o014.li2.chapter2.glimpse-of-lattice-theory.p006
Untuk memperagakan operasi $\wedge$ dan $\vee$, mari kita buktikan fakta
berikut. Jika $a^\flat,a,b \in P$ memenuhi $a^\flat \leq a$, maka
% segment-id: o014.li2.chapter2.glimpse-of-lattice-theory.q001
\begin{equation}\label{eqn:modular-inequality}
	a^\flat \vee (a \wedge b) \leq (a^\flat \vee b) \wedge a.
\end{equation}
Berdasarkan definisi supremum, pertama-tama masalahnya tereduksi menjadi
pembuktian $a^\flat \leq (a^\flat \vee b) \wedge a$ dan
$(a \wedge b) \leq (a^\flat \vee b) \wedge a$. Berdasarkan definisi infimum,
hal ini selanjutnya masing-masing tereduksi menjadi pembuktian
\[ a^\flat \leq a^\flat \vee b, \quad a^\flat \leq a, \quad a \wedge b \leq a^\flat \vee b, \quad a \wedge b \leq a; \]
ketaksamaan ketiga mengikuti dari $a \wedge b \leq b \leq a^\flat \vee b$,
sedangkan yang lain sudah jelas.

% segment-id: o014.li2.chapter2.glimpse-of-lattice-theory.d002
\begin{definisi}\label{def:modular-lattice}
	\index{kisi!modular (modular)}
	\index{komplemen (complement)}
	Misalkan $P$ merupakan kisi.
	% segment-id: o014.li2.chapter2.glimpse-of-lattice-theory.l001
	\begin{itemize}
		% segment-id: o014.li2.chapter2.glimpse-of-lattice-theory.i001
		\item $P$ disebut \emph{kisi modular} jika sifat berikut berlaku: apabila
		$a^\flat,a,b \in P$ memenuhi $a^\flat \leq a$, maka
		% segment-id: o014.li2.chapter2.glimpse-of-lattice-theory.q002
		\[ a^\flat \vee (a \wedge b) = (a^\flat \vee b) \wedge a. \]
		% segment-id: o014.li2.chapter2.glimpse-of-lattice-theory.i002
		\item Misalkan $P$ merupakan kisi terbatas. Jika $x,c \in P$ memenuhi
		$x \vee c=1$ dan $x \wedge c=0$, maka $c$ disebut \emph{komplemen} dari
		$x$ di $P$.
	\end{itemize}
\end{definisi}

% segment-id: o014.li2.chapter2.glimpse-of-lattice-theory.p007
Sekarang mari kita tinjau contoh-contoh dasar kisi.
% segment-id: o014.li2.chapter2.glimpse-of-lattice-theory.l002
\begin{itemize}
	% segment-id: o014.li2.chapter2.glimpse-of-lattice-theory.i003
	\item Himpunan bilangan bulat positif $\Z_{\geq 1}$ merupakan kisi terhadap
	relasi keterbagian ($x \mid y \iff x \leq y$):
	$x \vee y=\mathrm{lcm}(x,y)$ dan $x \wedge y=\mathrm{gcd}(x,y)$. Silakan
	pembaca memeriksa bahwa kisi ini modular, mempunyai batas bawah $1$, tetapi
	tidak mempunyai batas atas.

	% segment-id: o014.li2.chapter2.glimpse-of-lattice-theory.i004
	\item Ambil sebuah gelanggang $R$ dan modul-$R$ kiri $M$. Himpunan submodulnya
	$\mathrm{Sub}_M$ merupakan kisi modular terbatas terhadap $\subset$:
	$x \vee y=x+y$ dan $x \wedge y=x \cap y$; batas atasnya adalah $M$ dan batas
	bawahnya adalah $\{0\}$. Inilah asal-usul istilah ``kisi modular''.

	% segment-id: o014.li2.chapter2.glimpse-of-lattice-theory.i005
	\item Semua subruang tertutup dari ruang Hilbert $H$ merupakan kisi terbatas
	terhadap $\subset$: $x \vee y:=\overline{x+y}$ dan $x \wedge y=x \cap y$.
	Dapat dibuktikan bahwa kisi ini tidak modular jika $H$ berdimensi tak hingga
	(Birkhoff--von Neumann).
	% Rujukan yang mungkin: Proposition 4.4 dalam https://link.springer.com/chapter/10.1007/978-94-015-9026-6_4
\end{itemize}

% segment-id: o014.li2.chapter2.glimpse-of-lattice-theory.r001
\begin{catatan}\label{rem:lattice-duality}
	Definisi kisi juga mengandung suatu dualitas. Untuk urutan parsial
	$\leq$ yang diberikan pada himpunan $P$, kita dapat meninjau urutan parsial
	lawannya $\leq^{\opp}$: $x \leq^{\opp}y \iff x \geq y$; ini sama artinya
	dengan mengganti $\mathcal{P}$ dengan $\mathcal{P}^{\opp}$. Jika
	$(P,\leq)$ merupakan kisi (atau kisi terbatas), maka demikian pula
	$(P,\leq^{\opp})$. Lebih lanjut, $\wedge,\vee,0,1$ dalam $(P,\leq)$
	masing-masing bersesuaian dengan $\vee,\wedge,1,0$ dalam
	$(P,\leq^{\opp})$.
\end{catatan}

% segment-id: o014.li2.chapter2.glimpse-of-lattice-theory.le001
\begin{lema}\label{prop:modularity-criterion}
	Misalkan $P$ merupakan kisi. Maka $P$ merupakan kisi modular jika dan hanya
	jika sifat berikut berlaku untuk setiap interval $I$ dalam $P$: jika
	$c^\flat,c$ sama-sama merupakan komplemen dari $x \in I$ di $I$ dan
	$c^\flat \leq c$, maka $c^\flat=c$.
\end{lema}
% segment-id: o014.li2.chapter2.glimpse-of-lattice-theory.pf001
\begin{bukti}
	Jika $P$ merupakan kisi modular, maka setiap interval $I$ juga merupakan
	kisi modular. Tanpa mengurangi keumuman, anggap $I=P$; untuk
	$x,c^\flat,c$ dalam pernyataan tersebut, kita mempunyai
	% segment-id: o014.li2.chapter2.glimpse-of-lattice-theory.q003
	\[ c = c \wedge 1 = c \wedge (x \vee c^\flat) = c^\flat \vee (x \wedge c) = c^\flat \vee 0 = c^\flat . \]

	Sebaliknya, anggap syarat tentang komplemen berlaku untuk setiap interval.
	Misalkan $a^\flat,a,b \in P$ memenuhi $a^\flat \leq a$. Tetapkan
	% segment-id: o014.li2.chapter2.glimpse-of-lattice-theory.q004
	\[ b \wedge a \leq c_1 := a^\flat \vee (b \wedge a) \underset{\because \eqref{eqn:modular-inequality}}{\leq} a \wedge (b \vee a^\flat) =: c_2 \leq b \vee a^\flat. \]
	Kita akan membuktikan bahwa $c_1,c_2$ sama-sama merupakan komplemen dari
	$b$ dalam $\left[b \wedge a,\; b \vee a^\flat\right]$, sehingga
	$c_1=c_2$. Pertama, dari definisi langsung diperoleh $c_1 \leq a$; oleh
	karena itu,
	% segment-id: o014.li2.chapter2.glimpse-of-lattice-theory.q005
	\[ b \wedge a \geq b \wedge c_1 = (a^\flat \vee (b \wedge a)) \wedge b \underset{\because \eqref{eqn:modular-inequality}}{\geq} (a^\flat \wedge b) \vee (b \wedge a) = b \wedge a , \]
	sehingga $b \wedge c_1=b \wedge a$. Di sisi lain,
	$b \vee c_1=a^\flat \vee (b \wedge a) \vee b=a^\flat \vee b$. Jadi,
	$c_1$ memang merupakan komplemen dari $b$ dalam
	$\left[b \wedge a,\; b \vee a^\flat\right]$.

	Dualitas kisi (Catatan~\ref{rem:lattice-duality}) menukar hubungan urutan
	antara $a^\flat,a$, menukar $\wedge$ dan $\vee$, serta menukar peran $c_1$
	dan $c_2$. Dengan demikian, argumen di atas yang dijalankan dalam
	$(P,\leq^{\opp})$ memperlihatkan bahwa $c_2$ merupakan komplemen dari $b$
	dalam $\left[b \wedge a,\; b \vee a^\flat\right]$.
\end{bukti}

% segment-id: o014.li2.chapter2.glimpse-of-lattice-theory.pr001
\begin{proposisi}\label{prop:lattice-diamond-isom}
	\index{isomorfisme standar (standard isomorphism)}
	Misalkan $a,b$ merupakan unsur kisi modular $(P,\leq)$. Maka terdapat
	isomorfisme antarhimpunan terurut parsial
	% segment-id: o014.li2.chapter2.glimpse-of-lattice-theory.g001
	\[\begin{tikzcd}[row sep=tiny]
		{[a \wedge b, a]} \arrow[leftrightarrow, r, "1:1"] & {[b, a \vee b]} \\
		x \arrow[mapsto, r] & x \vee b \\
		y \wedge a & y \arrow[mapsto, l]
	\end{tikzcd}\]
	yang disebut \emph{isomorfisme standar}.
\end{proposisi}
% segment-id: o014.li2.chapter2.glimpse-of-lattice-theory.pf002
\begin{bukti}
	Kedua pemetaan jelas terdefinisi dengan baik dan mempertahankan urutan.
	Untuk $x \in [a \wedge b,a]$, definisi kisi modular memberikan
	$(x \vee b) \wedge a=(a \wedge b) \vee x$, dan ruas kanannya adalah $x$.
	Secara dual, $y \in [b,a \vee b]$ mengakibatkan
	$(y \wedge a) \vee b=y$ (Catatan~\ref{rem:lattice-duality}); jadi kedua
	pemetaan itu saling invers.
\end{bukti}

% segment-id: o014.li2.chapter2.glimpse-of-lattice-theory.p008
Lema Zassenhaus dalam aljabar elementer (lihat \cite[Lema 4.6.4]{Li1}, lebih
tepatnya versi teori modulnya) beserta konsekuensi-konsekuensinya dapat
disarikan dengan menggunakan teori kisi.

% segment-id: o014.li2.chapter2.glimpse-of-lattice-theory.th001
\begin{teorema}[Lema Zassenhaus untuk Kisi Modular]\label{prop:lattice-Zassenhaus}
	\index{Lema Zassenhaus (Zassenhaus Lemma)}
	Misalkan $(P,\leq)$ merupakan kisi modular. Untuk unsur-unsur
	$u^\flat \leq u$ dan $v^\flat \leq v$ yang diberikan, terdapat diagram
	berikut:
	% segment-id: o014.li2.chapter2.glimpse-of-lattice-theory.g002
	\[ \begin{tikzcd}[column sep=tiny, row sep=small]
		u^\flat \vee (u \wedge v) \arrow[dash, dd] \arrow[dash, rd] & & (u \wedge v) \vee v^\flat \arrow[dash, dd] \arrow[dash, ld] \\
		& u \wedge v \arrow[dash, dd] & \\
		u^\flat \vee (u \wedge v^\flat) \arrow[dash, rd] & & (u^\flat \wedge v) \vee v^\flat \arrow[dash, ld] \\
		& (u^\flat \wedge v) \vee (u \wedge v^\flat)
	\end{tikzcd} \]
	yang berarti bahwa unsur-unsurnya memenuhi pola-pola berikut:
	% segment-id: o014.li2.chapter2.glimpse-of-lattice-theory.q006
	\begin{equation}\label{eqn:Zassenhaus-cases}
	% segment-id: o014.li2.chapter2.glimpse-of-lattice-theory.g003
	\begin{tikzcd}
		x \arrow[dash, d, "\geq" description, sloped] \\ y
	\end{tikzcd} \qquad
	% segment-id: o014.li2.chapter2.glimpse-of-lattice-theory.g004
	\begin{tikzcd}[column sep=tiny]
		x \arrow[dash, rd] & & y \arrow[dash, ld] \\
		& x \wedge y &
	\end{tikzcd} \quad
	% segment-id: o014.li2.chapter2.glimpse-of-lattice-theory.g005
	\begin{tikzcd}[column sep=tiny]
		{} & x \vee y \arrow[dash, ld] \arrow[dash, rd] & \\
		x & & y
	\end{tikzcd} .
	\end{equation}
	Selain itu, terdapat isomorfisme standar antarinterval
	% segment-id: o014.li2.chapter2.glimpse-of-lattice-theory.q007
	\begin{multline*}
		\left[ u^\flat \vee (u \wedge v^\flat), \; u^\flat \vee (u \wedge v) \right] \leftiso \left[ (u^\flat \wedge v) \vee (u \wedge v^\flat), \; u \wedge v \right] \\
		\rightiso \left[ (u^\flat \wedge v) \vee v^\flat, \; (u \wedge v) \vee v^\flat \right],
	\end{multline*}
	yang direalisasikan oleh pemetaan
	$x \vee u^\flat \vee (u \wedge v^\flat) \mapsfrom x \mapsto
	x \vee (u^\flat \wedge v) \vee v^\flat$ dalam
	Proposisi~\ref{prop:lattice-diamond-isom}.
\end{teorema}
% segment-id: o014.li2.chapter2.glimpse-of-lattice-theory.pf003
\begin{bukti}
	Isomorfisme pada bagian terakhir diperoleh dengan menerapkan
	Proposisi~\ref{prop:lattice-diamond-isom} pada kedua jajaran genjang dalam
	diagram; pokoknya adalah memeriksa relasi-relasi dalam
	\eqref{eqn:Zassenhaus-cases}. Perhatikan bahwa diagram simetris kiri--kanan
	terhadap $u \leftrightarrow v$ dan $u^\flat \leftrightarrow v^\flat$;
	karena itu, cukup periksa bagian kirinya. Relasi urutan parsial di antara
	unsur-unsurnya langsung terlihat. Selanjutnya, untuk kasus $\wedge$,
	definisi kisi modular memberikan
	% segment-id: o014.li2.chapter2.glimpse-of-lattice-theory.q008
	\begin{equation*}
		\left( u^\flat \vee (u \wedge v^\flat)\right) \wedge (u \wedge v) = \left( u^\flat \wedge (u \wedge v) \right) \vee (u \wedge v^\flat) = (u^\flat \wedge v) \vee (u \wedge v^\flat).
	\end{equation*}
	Untuk kasus $\vee$, ketaksamaan $u \wedge v^\flat \leq u \wedge v$ langsung
	memberikan
	% segment-id: o014.li2.chapter2.glimpse-of-lattice-theory.q009
	\begin{equation*}
		u^\flat \vee (u \wedge v^\flat) \vee (u \wedge v) = u^\flat \vee (u \wedge v).
	\end{equation*}
	Dengan demikian, \eqref{eqn:Zassenhaus-cases} terbukti.
\end{bukti}

% segment-id: o014.li2.chapter2.glimpse-of-lattice-theory.p009
Selanjutnya, tinjau rantai menurun hingga dalam suatu kisi, yang ditulis
$x_0 \geq x_1 \geq \cdots \geq x_r$ ($r \in \Z_{\geq 0}$). Rantai menurun
yang diperoleh dengan menyisipkan berhingga banyak unsur perantara disebut
\emph{penghalusan} dari rantai menurun semula; jika unsur-unsur yang disisipkan
mencakup suatu $y \notin \{x_0,\ldots,x_r\}$, penghalusan itu disebut
penghalusan sejati.

% segment-id: o014.li2.chapter2.glimpse-of-lattice-theory.d003
\begin{definisi}\label{def:Schreier-equiv}
	Dua rantai menurun dengan panjang sama dalam suatu kisi modular,
	$x_0 \geq \cdots \geq x_r$ dan $x'_0 \geq \cdots \geq x'_r$, disebut
	\emph{ekuivalen} jika kondisi berikut dipenuhi: terdapat bijeksi
	$\sigma$ dari $\{0,\ldots,r-1\}$ ke dirinya sendiri (yakni suatu
	permutasi)\footnote{Catatan penerjemah (O014-C021): sumber menulis
	$\{0,\ldots,r\}$, padahal faktor-faktor berurutan yang dipasangkan berindeks
	$0\leq i<r$; indeks $i=r$ akan membuat $x_{i+1}$ tidak terdefinisi.}
	sehingga, untuk setiap $0\leq i<r$, terdapat isomorfisme antarhimpunan
	terurut parsial
	% segment-id: o014.li2.chapter2.glimpse-of-lattice-theory.q010
	\[ \tau_i: \left[ x_{i+1}, x_i \right] \simeq \left[ x'_{\sigma(i) + 1}, x'_{\sigma(i)} \right], \]
	dan $\tau_i$ terurai menjadi komposisi berhingga dari isomorfisme standar
	(Proposisi~\ref{prop:lattice-diamond-isom}) atau invers-inversnya.
\end{definisi}

% segment-id: o014.li2.chapter2.glimpse-of-lattice-theory.th002
\begin{teorema}[Teorema Penghalusan Schreier untuk Kisi Modular]\label{prop:Schreier-refinement}
	\index{Teorema Penghalusan Schreier (Schreier Refinement Theorem)}
	\mbox{}\par
	Tinjau rantai-rantai menurun
	% segment-id: o014.li2.chapter2.glimpse-of-lattice-theory.q011
	\[ x_0 \geq \cdots \geq x_r, \quad y_0 \geq \cdots \geq y_s, \]
	dalam kisi modular $(P,\leq)$ sedemikian sehingga $x_0=y_0$ dan $x_r=y_s$,
	dengan $r,s \in \Z_{\geq 0}$. Maka masing-masing rantai mempunyai suatu
	penghalusan, dan kedua penghalusan itu ekuivalen.
\end{teorema}
% segment-id: o014.li2.chapter2.glimpse-of-lattice-theory.pf004
\begin{bukti}
	Argumennya sama seperti dalam kasus grup \cite[Teorema 4.6.6]{Li1}; keduanya
	berdasarkan Teorema~\ref{prop:lattice-Zassenhaus}. Kita uraikan garis
	besarnya. Definisikan
	% segment-id: o014.li2.chapter2.glimpse-of-lattice-theory.q012
	\[ x_{i, j} := x_{i+1} \vee (x_i \wedge y_j), \quad y_{j, i} := (x_i \wedge y_j) \vee y_{j+1}, \]
	dengan $0 \leq i<r$ dan $0 \leq j \leq s$ untuk $x_{i,j}$, serta
	$0 \leq i \leq r$ dan $0 \leq j<s$ untuk $y_{j,i}$. Maka
	$x_{i,j+1} \leq x_{i,j}$, $x_{i,0}=x_i$, dan $x_{i,s}=x_{i+1}$, sehingga
	$\left(x_{i,j}\right)_{i,j}$, menurut urutan ini, memperhalus
	$x_0 \geq \cdots \geq x_r$. Demikian pula,
	$\left(y_{j,i}\right)_{j,i}$ memperhalus $y_0 \geq \cdots \geq y_s$.
	Dengan mengambil $u^\flat=x_{i+1}$, $u=x_i$, $v^\flat=y_{j+1}$, dan
	$v=y_j$ dalam Teorema~\ref{prop:lattice-Zassenhaus}, kita memperoleh
	% segment-id: o014.li2.chapter2.glimpse-of-lattice-theory.q013
	\[ \left[ x_{i, j+1}, x_{i,j} \right] \simeq \left[ y_{j, i+1}, y_{j, i} \right], \]
	dan isomorfisme ini dapat diuraikan menjadi isomorfisme standar beserta
	invers-inversnya. Terbuktilah pernyataan tersebut.
\end{bukti}

% segment-id: o014.li2.chapter2.glimpse-of-lattice-theory.p010
Mulai sekarang, kita hanya meninjau rantai naik atau rantai menurun yang ketat.

% segment-id: o014.li2.chapter2.glimpse-of-lattice-theory.d004
\begin{definisi}\index{deret komposisi (composition series)}
	Tetapkan himpunan terurut parsial $P$ dan unsur-unsur $a\leq b$ di
	dalamnya.\footnote{Catatan penerjemah (O014-C022): sumber menulis $a<b$,
	tetapi kalimat berikutnya secara eksplisit mencakup $a=b$ dan $r=0$; target
	menggunakan hipotesis lemah yang konsisten dengan kasus tersebut.} Jika
	rantai menurun $b=x_0>\cdots>x_r=a$ dalam $P$ tidak mempunyai penghalusan
	sejati, rantai itu disebut \emph{deret komposisi} dari $[a,b]$.
\end{definisi}

% segment-id: o014.li2.chapter2.glimpse-of-lattice-theory.p011
Panjang deret komposisi tersebut didefinisikan sebagai $r \in \Z_{\geq 0}$;
perhatikan bahwa $r=0$ jika dan hanya jika $a=b$.
\index{panjang (length)}

% segment-id: o014.li2.chapter2.glimpse-of-lattice-theory.th003
\begin{teorema}[Teorema Jordan--Hölder untuk Kisi Modular]\label{prop:lattice-JH}
	\index{Jordan-Holder@Teorema Jordan--Hölder (Jordan--Hölder Theorem)}
	Misalkan $a<b$ merupakan unsur kisi modular $(P,\leq)$. Maka setiap dua
	deret komposisi dari $[a,b]$ mempunyai panjang yang sama dan saling
	ekuivalen dalam pengertian Definisi~\ref{def:Schreier-equiv}.
\end{teorema}
% segment-id: o014.li2.chapter2.glimpse-of-lattice-theory.pf005
\begin{bukti}
	Pernyataan ini merupakan akibat langsung dari
	Teorema~\ref{prop:Schreier-refinement}.
\end{bukti}

% segment-id: o014.li2.chapter2.glimpse-of-lattice-theory.p012
Persoalan pentingnya ialah menentukan interval dalam himpunan terurut parsial
mana yang mempunyai deret komposisi. Seperti pada situasi teori modul yang
sudah dikenal, keberadaan tersebut dapat dijamin oleh syarat rantai naik dan
rantai menurun.

% segment-id: o014.li2.chapter2.glimpse-of-lattice-theory.d005
\begin{definisi}\label{def:lattice-Noether-Artin}
	\index{himpunan terurut parsial!berpanjang hingga (finite length)}
	\index{himpunan terurut parsial!Artinian, Noetherian}
	Misalkan $(P,\leq)$ merupakan himpunan terurut parsial. Jika di $P$ tidak
	ada rantai naik tak hingga $x_1<x_2<x_3<\cdots$, maka $P$ disebut
	\emph{Noetherian}; jika tidak ada rantai menurun tak hingga
	$x_1>x_2>x_3>\cdots$, maka $P$ disebut \emph{Artinian}. Himpunan terurut
	parsial yang sekaligus Artinian dan Noetherian disebut
	\emph{berpanjang hingga}.
\end{definisi}

% segment-id: o014.li2.chapter2.glimpse-of-lattice-theory.p013
Jika $(P,\leq)$ bersifat Noetherian (atau Artinian, atau berpanjang hingga),
maka setiap himpunan bagiannya juga demikian. Mudah dibuktikan bahwa syarat
Noetherian (atau Artinian) ekuivalen dengan pernyataan bahwa setiap himpunan
bagian tak kosong dari $P$ mempunyai unsur maksimal (atau unsur minimal)
terhadap $\leq$. Dalam bab ini, konsep-konsep tersebut terutama akan diterapkan
pada himpunan terurut parsial dari subobjek; lihat
Definisi~\ref{def:sub-quot-obj}.

% segment-id: o014.li2.chapter2.glimpse-of-lattice-theory.d006
\begin{definisi}\label{def:finite-length-object}
	\index{objek!berpanjang hingga (of finite length)}
	\index{objek!Noetherian (Noetherian)}
	\index{objek!Artinian (Artinian)}
	Tetapkan suatu kategori $\mathcal{C}$. Objek $X$ dari $\mathcal{C}$ disebut
	Noetherian (atau Artinian, atau berpanjang hingga) jika himpunan terurut
	parsial subobjeknya $(\mathrm{Sub}_X,\subset)$ bersifat Noetherian (atau
	Artinian, atau berpanjang hingga).
\end{definisi}

% segment-id: o014.li2.chapter2.glimpse-of-lattice-theory.p014
Sekarang kita kembali ke himpunan terurut parsial umum.

% segment-id: o014.li2.chapter2.glimpse-of-lattice-theory.le002
\begin{lema}\label{prop:lattice-composition-series}
	Tinjau unsur-unsur $a \leq b$ dari himpunan terurut parsial $(P,\leq)$.
	% segment-id: o014.li2.chapter2.glimpse-of-lattice-theory.l003
	\begin{enumerate}[(i)]
		% segment-id: o014.li2.chapter2.glimpse-of-lattice-theory.i006
		\item Jika $[a,b]$ berpanjang hingga, maka setiap rantai di dalamnya dapat
		diperhalus menjadi suatu deret komposisi; khususnya, $[a,b]$ mempunyai
		deret komposisi.
		% segment-id: o014.li2.chapter2.glimpse-of-lattice-theory.i007
		\item Jika $[a,b]$ diasumsikan merupakan kisi modular dan mempunyai deret
		komposisi, maka $[a,b]$ berpanjang hingga.
	\end{enumerate}
\end{lema}
% segment-id: o014.li2.chapter2.glimpse-of-lattice-theory.pf006
\begin{bukti}
	Tanpa mengurangi keumuman, anggap $a<b$. Untuk (i), jika diberikan rantai
	$y_0>\cdots>y_r$, cukup ditunjukkan bahwa setiap $[y_{i+1},y_i]$ mempunyai
	deret komposisi.\footnote{Catatan penerjemah (O014-C023): sumber menulis
	$[y_i,y_{i+1}]$ dan kemudian menetapkan $y_i=a$, $y_{i+1}=b$. Karena
	$y_i>y_{i+1}$, kedua urutan itu dibalik; target menggunakan interval
	$[y_{i+1},y_i]$ dan penetapan yang bertipe benar.} Tanpa mengurangi
	keumuman, anggap $y_{i+1}=a$ dan $y_i=b$. Syarat Artinian menjamin adanya
	unsur minimal
	$x^0 \in [a,b]$ dengan $x^0>a$. Demikian pula, jika $x^0 \neq b$, terdapat
	unsur minimal $x^1 \in [a,b]$ dengan $x^1>x^0$, dan seterusnya. Berdasarkan
	syarat Noetherian, proses ini mesti berhenti setelah berhingga banyak langkah,
	dan menghasilkan rantai $b>\cdots>x^0>a$ tanpa penghalusan sejati.

	Untuk (ii), pilih suatu deret komposisi $b=x_0>\cdots>x_r=a$.
	Teorema~\ref{prop:Schreier-refinement} menunjukkan bahwa setiap rantai
	hingga $\cdots>y_i>y_{i+1}>\cdots$ mempunyai penghalusan yang
	ekuivalen dengan deret komposisi di atas; karena itu, panjang rantai tersebut
	tidak dapat melampaui panjang deret komposisi. Jadi, $[a,b]$ merupakan
	himpunan terurut parsial berpanjang hingga.
\end{bukti}

% segment-id: o014.li2.chapter2.glimpse-of-lattice-theory.d007
\begin{definisi}\label{def:lattice-length}
	Misalkan $(P,\leq)$ merupakan kisi modular terbatas yang berpanjang hingga.
	Pilih sebarang deret komposisi $1=x_0>\cdots>x_r=0$ dari $P$. Bilangan
	$r \in \Z_{\geq 0}$ disebut \emph{panjang} dari $(P,\leq)$.
\end{definisi}

% segment-id: o014.li2.chapter2.glimpse-of-lattice-theory.p015
Teorema~\ref{prop:lattice-JH} menunjukkan bahwa panjang tersebut terdefinisi
dengan baik dan tidak bergantung pada deret komposisi yang dipilih; panjang
$(P,\leq)$ adalah $0$ jika dan hanya jika $P$ merupakan himpunan satu titik.
